\section{Background and construction}
\label{sec:background}

\subsection{Notation}

Let \(\F\) be the Mersenne prime field used for trace values, and let \(\K\) be its
degree-four extension used for verifier challenges.  We write \(H\) for a
domain-separated hash function and use distinct domains for note commitments,
nullifiers, Merkle nodes, and transcript operations.  A \emph{transparent}
proof uses public hash-based commitments and therefore requires no structured
reference string or ceremony.

The proof system follows the interactive-oracle-proof view underlying STARKs
\cite{BenSassonChiesaSpooner2016,
BenSassonBentovHoreshRiabzev2018STARK}.  The prover commits to encoded
execution data, the verifier derives public challenges from the transcript,
and a FRI-style test checks proximity to a low-degree codeword
\cite{BenSassonBentovHoreshRiabzev2018FRI}.  Here \emph{FRI} refers to the
low-degree test: the verifier queries a few committed positions and repeatedly
folds the alleged codeword until only a small polynomial remains.

\subsection{Private-spend relation}

The public statement is
\[
  x=(P,s,R_{\mathrm{old}},N,C_{\mathrm{out}},R_{\mathrm{new}},a,f,d),
\]
where \(P\) identifies the pool, \(s\) is its sequence number,
\(R_{\mathrm{old}}\) and \(R_{\mathrm{new}}\) are the old and new note-tree
roots, \(N\) is the input nullifier, \(C_{\mathrm{out}}\) is the output-note
commitment, \(a\) is the asset identifier, \(f\) is the public fee, and \(d\)
binds the statement to a deployment domain.  A nullifier is a public,
one-time tag derived from the input note's secret opening.  Publishing it
prevents a second successful spend without revealing which commitment in the
note tree was consumed.

The private witness contains the input-note opening and ownership credential,
the input note's depth-twenty Merkle path, and the opening data for the output
note.  With input and output values \(v_{\mathrm{in}}\) and
\(v_{\mathrm{out}}\), the relation enforces
\begin{equation}
  v_{\mathrm{in}}=v_{\mathrm{out}}+f,
  \qquad
  0\leq v_{\mathrm{in}},v_{\mathrm{out}},f<2^{30}.
  \label{eq:value-balance}
\end{equation}
The range conditions make the field equation an integer conservation law,
ruling out modular wraparound.  Both notes use asset \(a\).  The input
credential determines its owner and nullifier, while each note commitment
binds the owner, value, asset, and note randomness.

The input path reconstructs \(R_{\mathrm{old}}\).  Replacing its leaf by the
output commitment at the same depth-twenty position, with the same direction
bits and siblings, reconstructs \(R_{\mathrm{new}}\).  Thus the evaluated
statement is a one-input, one-output replacement in a fixed-position note
tree.

Note commitments, nullifiers, and trace Merkle nodes use Poseidon2, an
arithmetic-oriented hash function designed for efficient use inside proof
systems \cite{GrassiKhovratovichSchofnegger2023Poseidon2}.

\begin{definition}[Valid private spend]
\label{def:spend}
A public statement \(x\) is valid if there exists a private witness satisfying
the range and value-conservation conditions, ownership and nullifier
derivations, input and output note-commitment equations, and both Merkle-root
equations described above.
\end{definition}

\subsection{Evaluated proof parameters}

The prover commits to sixteen base-field columns and three extension-field
columns.  A nonzero challenge \(\gamma\in\K\) combines these nineteen columns as
\(\sum_{i=0}^{18}\gamma^i c_i\).  The initial codeword is evaluated on a
circle domain over \(\F\); four circle symbols form one query fibre.  Four
radix-four folds reduce it through three later line domains to a polynomial
with four coefficients.  Five salted Merkle trees authenticate the two
initial commitments and the three later layers.

\begin{table}[t]
\centering
\small
\caption{Parameters of the evaluated proof.}
\label{tab:profile-parameters}
\begin{tabular}{@{}ll@{}}
\toprule
Quantity & Value \\
\midrule
Base field & \(\mathbb F_{2^{31}-1}\) \\
Challenge field & degree-four extension \(\K\) \\
Trace rows & \(2^{10}\) \\
Initial evaluation symbols & \(2^{19}\) \\
Initial query fibres & \(2^{17}=131{,}072\) \\
Batched columns & 19 \\
Distinct queries & 18, selected from at most 64 draws \\
Low-degree folds & four radix-four rounds \\
Final polynomial & four coefficients in \(\K\) \\
Proof-of-work difficulties & 37, 34, 33, 30, 25, and 32 bits \\
Public-coin boundaries & 30 after the initial commitment \\
Authenticated trees & five \\
Proof body & 75,358 bytes \\
\bottomrule
\end{tabular}
\end{table}

The circle-domain construction follows work on circle codes and Circle
STARKs \cite{HaboeckLubarovNabaglo2023,HaboeckLevitPapini2024,
CarmonEtAl2026Stwo}.  The evaluated profile uses its own four-way folding
schedule, relation messages, bounded distinct-query sampler, and sequential
work schedule; \cref{sec:soundness} gives the corresponding checked reduction.

\subsection{Transcript and query selection}

The Fiat--Shamir transform replaces the verifier's random challenges with
hashes of the public transcript.  The transcript absorbs the proof parameters and
statement, commitment roots and public salts, algebraic claims, relation
points and messages, later-layer roots, final polynomial, and six work
nonces.  Each nonce is tested against the current transcript state before
being absorbed; the protected challenge is sampled only afterwards.  Typed
labels and length-delimited payloads separate the operations.

A work nonce succeeds when the corresponding SHA-256 digest has the required
number of leading zero bits.  Finding it takes repeated trials.  The
random-oracle model used in the security reduction treats the hash output on
each new input as an independent uniform value while returning the same value
on repeated inputs.

After every commitment, relation message, fold challenge, nonce, and final
coefficient has been fixed, the verifier derives query material from SHA-256.
It interprets each 32-byte block as eight little-endian 32-bit words, masks
each word into the \(2^{17}\)-element query domain, and keeps first
occurrences.  Verification proceeds only if eighteen distinct positions are
obtained within sixty-four draws.

\subsection{Atomic state transition}

The proof body is sealed in a program-owned account before the spend
transaction.  The spend instruction validates the public statement and proof,
checks that the nullifier has no existing marker, and then records both the
marker and the new pool state in the same Solana transaction.  Solana's
transaction semantics make these writes atomic when the instruction succeeds.
A later cleanup transaction may close the retained proof account and refund
its balance; the nullifier and pool update remain in state.
